\scanpage{62}\mspage{61}
\deleted{Soit} Ainsi, les \unsure{systèmes|système} $(\mathcal F,\leq)$ satisfaisant $C1,C2,C3$ et tels que toute $F\in\mathcal F$ soit ``finie'' (i.e. l'ens $\mathcal F_{\leq F}$ des sous-figures est fini, ou encore, l'ens $M_{\leq F}\inserted{=\overline F}$ \deleted{des} de \deleted{sous} ses strates est fini) équivaut à :

$(M,\leq,R)$ \deleted{\unsure{str.}} où $\leq$ relation d'ordre, $R$ relation refl. et sym., satisfais. à
\[
\begin{aligned}
\text{a) }&\text{les }M_{\leq x}\text{ finis}\quad(x\in M)\\
\text{b) }&\text{Si }(x,y)\in R\text{ et }x'\leq x,y'\leq y,\ \text{on a }\deleted{R}(x',y')\in R\,;
\end{aligned}
\]
où les figures sont interprétées comme les parties \uline{fermées}, \uline{finies} \inserted{$A$} de $M$, telles que $\forall x,y\in A$, on ait $(x,y)\in R$.

% slanted
\margin{\unsure{Ainsi} la notion \uline{\unsure{combinatoire}} de connexité. Elle \illegible{1} en évidence \deleted{que} aussi la notion \illegible{1} qui la sous-entend, que les strates $x$\unsure{,} sont 0-connexes + \ldots\ En fait, si on \unsure{veut|voulait} \deleted{``exprimer''} \unsure{$=$|--} que les \unsure{conditions} pr. $A$ clos \deleted{\unsure{et}} de $x$ \unsure{soient|soit} \unsure{co}}

\uline{Composantes connexes} Une figure est dite \uline{connexe} si \deleted{\unsure{on a}}
\[
F=F'\cup F''\quad
\inserted{\text{avec }F',F''\ \text{\uline{disjoints} i.e. }F'\cap F''=\varnothing}
\ \Longrightarrow\ F'=\deleted{\unsure{\varnothing}}F\ \text{ou }F''=\deleted{\unsure{\varnothing}}F.
\]
\margin{\uline{NB} $F$ connexe ss $\overline F$ l'est comme ens. ordonné}

% slanted
\margin{On écrit $F=F'\amalg F''$ (signe de la somme disjointe\unsure{.|)} (\uline{NB} Si \deleted{deux} parmi les \deleted{fig.} $F',F'',F$ sont connexes, la troisième aussi.}

Une figure est dite \uline{0-connexe} si elle est connexe et $\ne\varnothing_{\mathcal F}$. Une sous-figure $F_0$ de $F$ est appelée \uline{composante connexe} de $F$ si a) elle est 0-connexe et

b) \deleted{$F_0$ est une \unsure{compos}} \inserted{$F=F_0\amalg G$, $G$ compatible} \deleted{avec $F$, disjoint de $F_0$.}

\deleted{\unsure{Il est}} \deleted{\uline{NB}}\inserted{Lemme} Soit $F=F'\amalg F''$, \deleted{avec $F'\cap F''=\varnothing$}, et soit $G$ une sous-figure connexe de $G$. Alors $G$ est contenue ds $F'$ ou $F''$.
